\newpage
\phantomsection
\label{prf:zero-product}

\begin{remark*}[Return]
\hyperref[thm:zero-product]{Return to Theorem}
\end{remark*}

\begin{theorem*}[Zero Product]
Let $F$ be a field and let $a\in F$.  Then
\[
a\cdot 0 = 0.
\]
\end{theorem*}

\begin{proof}
\textbf{Professional Standard Proof.}~\\
Since $0$ is the additive identity, $0 + 0 = 0$.  By distributivity,
\[
a\cdot 0 = a\cdot(0 + 0) = a\cdot 0 + a\cdot 0.
\]
Adding $-(a\cdot 0)$ to both sides,
\[
0 = a\cdot 0.
\]
\end{proof}

\begin{proof}
\textbf{Detailed Learning Proof.}~\\

\textbf{Step 1.} Use the additive identity law to write $0$ as $0 + 0$.

By the field axiom for the additive identity,
\[
0 + 0 = 0.
\]

\textbf{Step 2.} Multiply both sides on the left by $a$.

\[
a \cdot 0 = a \cdot (0 + 0).
\]

\textbf{Step 3.} Apply distributivity.

By the distributivity axiom,
\[
a \cdot (0 + 0) = a\cdot 0 + a\cdot 0.
\]
Combining Steps 2 and 3:
\[
a\cdot 0 = a\cdot 0 + a\cdot 0.
\]

\textbf{Step 4.} Add the additive inverse of $a\cdot 0$ to both sides.

Adding $-(a\cdot 0)$ to both sides of $a\cdot 0 = a\cdot 0 + a\cdot 0$:
\[
a\cdot 0 + (-(a\cdot 0)) = a\cdot 0 + a\cdot 0 + (-(a\cdot 0)).
\]

\textbf{Step 5.} Apply the additive inverse axiom and associativity.

The left-hand side is $0$ by the additive inverse axiom.  On the right,
associativity of addition gives
\[
a\cdot 0 + (a\cdot 0 + (-(a\cdot 0))) = a\cdot 0 + 0 = a\cdot 0.
\]
Therefore
\[
0 = a\cdot 0.
\]

\textbf{Step 6.} Conclude.

By commutativity of equality, $a\cdot 0 = 0$.
\end{proof}

\begin{remark*}[Proof structure]
The key move is writing $0 = 0 + 0$ using the additive identity law, then
distributing $a$ across the sum to obtain $a\cdot 0 = a\cdot 0 + a\cdot 0$.
This says $a\cdot 0$ is a solution to $x = x + a\cdot 0$, and the only
solution to that equation is $x = 0$.  Cancellation is performed directly
by adding the additive inverse rather than invoking the cancellation theorem,
keeping the proof self-contained within the field axioms alone.
\end{remark*}

\begin{remark*}[Dependencies]~\\
\begin{itemize}
  \item \hyperref[def:field]{Definition~\ref*{def:field}}
        \textnormal{(additive identity: $0+0=0$; distributivity; additive inverse;
        associativity and commutativity of addition)}
\end{itemize}
\end{remark*}
